\documentclass[a4paper]{report}
% Shared preamble for the English probability and statistics documents.
% Based on the original course preamble. Compile with pdfLaTeX.
\usepackage[top=1.5cm,bottom=2cm,left=2cm,right=2cm]{geometry}
\usepackage[T1]{fontenc}
\usepackage{lmodern}
\usepackage[utf8]{inputenc}
\usepackage[english]{babel}

\usepackage{amsmath}
\usepackage{amssymb}
\usepackage{amsthm}
\usepackage{mathrsfs}
\usepackage{mathtools}
\usepackage{etoolbox}
\usepackage{tikz}
\usetikzlibrary{shapes.multipart,arrows.meta}
\usepackage{xcolor}
\usepackage{enumerate}
\usepackage[shortlabels]{enumitem}
\usepackage{hyperref}
\usepackage{version}
\usepackage{bbold}
\usepackage{fontawesome5}
\usepackage{fancyhdr}
\lfoot{}
\usepackage{natbib}
\usepackage{graphicx}
\usepackage{pdfpages}
\usepackage{authblk}
    
\usepackage{blindtext}
\usepackage{multicol}
\usepackage{collcell}
\usepackage{wrapfig}
\usepackage{blkarray}
\usepackage{cancel}
\usepackage{appendix}
\usepackage{float}
\makeatletter
\input{tex/assets/slashbox.sty}
\makeatother
\hypersetup{unicode,colorlinks,citecolor=red,filecolor=black,
  linkcolor=red,urlcolor=cyan}
\setlength{\columnsep}{1cm}
\pagenumbering{gobble}

\def\numtext#1{%
  \ifcase#1\or one\or two\or three\or four\or five\or six\or seven\or eight\or nine\or ten\or
  eleven\or twelve\or thirteen\or fourteen\or fifteen\or sixteen\or seventeen\or eighteen\or nineteen\or twenty\or Lots\fi}

\newcommand{\Byte}[1]{%
  \def\nodecontents{}%
  \foreach[count=\l] \k in {#1} {\xappto{\nodecontents}{\noexpand\nodepart{\numtext{\l}}\k}}%
  \node[rectangle split,rectangle split parts=8,rectangle split horizontal,draw] at (2,2){\nodecontents};}

\newcommand{\D}{{\mathrm{d}}}
\newcommand{\cE}{{\mathcal{E}}}
\newcommand{\cG}{{\mathcal{G}}}
\newcommand{\cX}{{\mathcal{X}}}
\newcommand{\cY}{{\mathcal{Y}}}
\newcommand{\cB}{{\mathcal{B}}}
\newcommand{\cF}{{\mathcal{F}}}
\newcommand{\cO}{{\mathcal{O}}}
\newcommand{\cS}{{\mathcal{S}}}
\newcommand{\cT}{{\mathcal{T}}}
\newcommand{\given}{\, \big| \,}
\newcommand{\cL}{{\mathscr{L}}}
\newcommand{\cD}{{\mathcal{D}}}
\newcommand{\cA}{{\mathscr{A}}}
\newcommand{\cP}{{\mathscr{P}}}
\newcommand{\cR}{{\mathscr{R}}}
\newcommand{\A}{{\mathcal A}}
\newcommand{\one}{\mathbb{1}}
\newcommand{\mathbbm}[1]{\mathbb{#1}}

\providecommand{\C}{}
\renewcommand{\C}{\mathbb C}
\newcommand{\R}{\mathbb R}
\newcommand{\Q}{\mathbb Q}
\newcommand{\N}{\mathbb N}
\newcommand{\T}{\mathbb T}
\newcommand{\F}{\mathbb F}
\renewcommand{\P}{\mathbb P}
\newcommand{\Z}{\mathbb Z}
\newcommand{\E}{\mathbb {E}}
\newcommand{\Pbb}[1]{\P\left( #1\right)}
\newcommand{\Ebb}[1]{\E\left[ #1\right]}

\newcommand{\lb}{\left (}
\newcommand{\rb}{\right )}
\newcommand{\lbb}{\left [}
\newcommand{\rbb}{\right ]}
\newcommand{\labs}{\left |}
\newcommand{\rabs}{\right |}
\newcommand{\lbrb}[1]{\lb #1 \rb}
\newcommand{\lbbrbb}[1]{\lbb#1\rbb}
\newcommand{\labsrabs}[1]{\labs#1\rabs}
\newcommand{\lbbrb}[1]{\lbb#1\rb}
\newcommand{\lbrbb}[1]{\lb#1\rbb}
\newcommand{\langlerangle}[1]{\langle#1\rangle}
\newcommand{\lbcurly}{\{}
\newcommand{\rbcurly}{\}}
\newcommand{\curly}[1]{\lbcurly#1\rbcurly}
\newcommand{\abs}{\labsrabs}
\newcommand{\ind}[1]{1_{#1}}
\newcommand{\indep}{\perp \!\!\! \perp}
\renewcommand{\d}{\, \mathrm{d}}

\newtheorem{theorem}{Theorem}
\newtheorem{problemCnt}{Lemma}
\theoremstyle{definition}
\newtheorem{problem}[problemCnt]{Problem}
\newtheorem{problembook}[problemCnt]{\faBook \,  Problem}
\newtheorem{problempencil}[problemCnt]{\faPencil* \,  Problem}
\newtheorem{problemstar}[problemCnt]{\faStar\,  Problem}
\newtheorem{problemstarpencil}[problemCnt]{\faStar \, \faPencil* \,   Problem}
\newtheorem{probleminteresting}[problemCnt]{\faSmile[regular] \,   Problem}
\newtheorem{solution}{Solution}

\newcommand{\fwcell}[1]{\makebox[\arraystretch\normalbaselineskip]{$#1$}}
\newenvironment{enumeratei}{\begin{enumerate}[(i)]}{\end{enumerate}}
\newenvironment{enumeratea}{\begin{enumerate}[(a)]}{\end{enumerate}}
\newenvironment{enumeraten}{\begin{enumerate}[\upshape 1.]}{\end{enumerate}}

% Compatibility helpers needed by a few files in the combined archive.
\AtBeginDocument{\noindent{\footnotesize English translation prepared with ChatGPT (OpenAI, GPT-6), 26 October 2026, for translation of the Bulgarian text into English and LaTeX formatting.
The Bulgarian originals remain the reference; translation may contain errors.}\par\smallskip}
\newcommand{\MissingGraphic}[1]{%
  \begin{center}\fbox{\parbox{0.82\linewidth}{\centering Missing source image: \texttt{#1}}}\end{center}}



\begin{document}
\noindent Applied Mathematics / Mathematics, 12 April 2025 \hfill Test 1
\hrulefill
\begin{it}The exact grading formula will depend on the results and will be no lower than 2 + the number of points. Time allowed: 4 hours. Good luck!\end{it}
\begin{problem} (0.5 points)
A box contains balls numbered $1,2,\ldots,14$. Three balls are drawn in succession without replacement. Find the probabilities of:
\begin{itemize}
\item $A=\{\text{all three numbers are even}\}$;
\item $B=\{\text{the largest number is 10}\}$;
\item $C=\{\text{the product is divisible by 3}\}$;
\item $D=\{\text{the numbers are consecutive, regardless of drawing order}\}$.
\end{itemize}
Are $A,B$ independent? What is the probability of $C$ given $A$?
\end{problem}
\begin{problem} (0.5 points)
The joint distribution of $X,Y$ is
\[\begin{array}{c|ccc}Y\backslash X&1&2&4\\\hline -1&3/15&1/15&0\\0&\alpha&2\alpha&3/15\\1&1/15&0&1/15\end{array}\]
Find $\alpha$ and $\mathrm{Cor}(X,Y)$.
\end{problem}
\begin{problem}
Let $X_1,\ldots,X_n$ be independent and identically distributed, and $S_n=X_1+\cdots+X_n$.
\begin{enumerate}
\item (0.25 points) Find $\E[X_i\mid S_n]$ for $1\leq i\leq n$. Recall that $\E[X\mid Y]$ takes the value $\E[X\mid Y=k]$ with probability $\P(Y=k)$.

For the following parts, $X_1\sim\mathrm{Geo}(p)$, meaning $\P(X=k)=(1-p)^{k-1}p$ for $k\geq1$.
\item (0.25 points) Find the probability generating function $g_{S_n}(s)$.
\item (0.25 points) Show that the geometric distribution is \emph{memoryless}: $\P(X>t+s\mid X>s)=\P(X>t)$. If an $\mathbb N$-valued random variable has this property, must it be geometric?
\item (0.25 points) Calculate $\P(X_1>X_2+X_3)$ and $\P(X_1>X_2+X_3\mid S_3=15)$.
\end{enumerate}
\end{problem}
\begin{problem}
The cells of a $9\times9$ grid are coloured randomly white, green or red (shown as white, grey and black below).
\[\includegraphics[width=0.4\textwidth]{tex/assets/prob-stat/grid2.png}\]
\begin{enumerate}
\item (0.25 points) What is the expected number of white or green cells? What is the probability that there are at least 60?

A horizontal (vertical) flag consists of three consecutive horizontal (vertical) cells coloured white, green and red from left to right (top to bottom).
\item (0.25 points) Is the number of horizontal flags binomially distributed? If so, with what parameters?
\item (0.25 points) What is the expected number of horizontal flags?
\item (0.25 points) Are the counts of horizontal and vertical flags independent? What is the expected total number of horizontal AND vertical flags?
\end{enumerate}
\end{problem}
\begin{problem}
Recall that $A$ occurs \emph{almost surely} if $\P(A)=1$. This is analogous to ``almost everywhere'' in analysis: $\P(A^c)=0$. We will study the Borel--Cantelli lemma, which informally states the following for a sequence $\mathcal A=(A_1,A_2,\ldots)$:
\begin{itemize}
\item If $\P(A_n)$ tends to 0 sufficiently fast, almost surely only finitely many events in $\mathcal A$ occur.
\item Otherwise, if the events are also independent, almost surely infinitely many occur.
\end{itemize}
Define $\displaystyle\limsup A_n=\bigcap_{n=1}^\infty\bigcup_{k=n}^\infty A_k$. What condition must $\omega\in\Omega$ satisfy to belong to $\limsup A_n$? Explain the connection with infinitely many events occurring.

\emph{Bonus (0.5 points):}
\begin{itemize}
\item Find $\limsup A_n$ for $A_n=[-1/n,3+1/n]$; for $A_n=[-2-(-1)^n,2+(-1)^n]$; and for $A_n=\{mp_n:m\in\mathbb N,\ p_k\text{ is the }k\text{th prime}\}$.
\item Compare $\limsup(A_n\cup B_n)$ and $\limsup(A_n\cap B_n)$ with $\limsup A_n$ and $\limsup B_n$.
\end{itemize}
Prove the following (or use a statement without proof in a subsequent part):
\begin{enumerate}
\item (0.25 points) If $B_1\supseteq B_2\supseteq\cdots$, then $\P(\bigcap_{n=1}^\infty B_n)=\lim_{n\to\infty}\P(B_n)$.
\item (0.25 points) \emph{Borel--Cantelli 1:} If $\sum_{n\geq1}\P(A_n)<\infty$, then $\P(\limsup A_n)=0$.
\item (0.25 points) $\P(\limsup A_n)=1-\lim_{N\to\infty}\P(\bigcap_{n\geq N}A_n^c)$.
\item (0.25 points) \emph{Borel--Cantelli 2:} If $\sum_{n\geq1}\P(A_n)=\infty$ and the events are independent, then $\P(\limsup A_n)=1$.
\end{enumerate}
\end{problem}
\begin{problem}
Angel tosses a fair coin infinitely often. Let $L_n$ be the longest run of consecutive heads among the first $n$ tosses. We seek the behaviour of $L_n$ as $n\to\infty$. For example:
\[\begin{array}{c|cccccccccc}n&1&2&3&4&5&6&7&8&9&\cdots\\\hline\text{Outcome}&T&H&H&T&H&H&H&T&H&\cdots\\L_n&0&1&2&2&2&2&3&3&3&\cdots\end{array}\]
Formally, let $X_i$ be iid $\mathrm{Ber}(1/2)$ variables and
\[L_n=\max\{k\geq1:\text{there exists }i\in\{1,\ldots,n-k+1\}\text{ such that }X_i=\cdots=X_{i+k-1}=1\}.\]
\begin{enumerate}
\item Prove for every integer $j\in[1,n]$ that
\[\P(L_n\geq j)\leq\frac{n-j+1}{2^j}\leq\frac n{2^j},\qquad\P(L_n\leq j)\leq(1-2^{-j})^{\lfloor n/j\rfloor}.\]
\item Find positive real numbers $a_n$ such that for every $\epsilon\in(0,1)$,
$\P(1-\epsilon<L_n/a_n<1+\epsilon)\to1$.
\item Prove for every $\epsilon\in(0,1)$ that
\[\P(L_n/a_n<1+\epsilon\text{ for all sufficiently large }n)=1,\]
\[\P(L_n/a_n>1-\epsilon\text{ for all sufficiently large }n)=1.\]
\item Conclude that $\P(L_n/a_n\to1)=1$.
\end{enumerate}
\end{problem}
\end{document}
